% META: {"id": "1SPE-SECDEG-EX-030", "chapitre": "1SPE-SECOND-DEGRE", "type_objet": "exercice", "capacites": ["1SPE-SECOND-DEGRE-C3", "1SPE-SECOND-DEGRE-C5"], "capacites_codes": ["C3", "C5"], "methodes": ["M3", "M5"], "parcours": 2, "competences": ["calculer", "raisonner", "modeliser"], "duree_min": 20, "mode_creation": "ex_nihilo", "sources_inspiration": [], "parametres_sympy": {}, "fichier_tex": "chapitres/1SPE-SECOND-DEGRE/exercices/1SPE-SECDEG-EX-030.tex", "corrige_tex": "chapitres/1SPE-SECOND-DEGRE/corriges/1SPE-SECDEG-CO-030.tex", "status": "generated"}
% BEGIN-VERIFY
% from sympy import symbols, solve, Rational, sqrt
% x = symbols('x')
% # h(t) = -t^2 + 6t + 7, t >= 0
% t = symbols('t')
% h = -t**2 + 6*t + 7
% # C3 : h(t) = 0 : delta = 36 + 28 = 64, t1 = (6+8)/2 = 7, t2 = (6-8)/2 = -1
% delta_h = 36 + 28
% assert delta_h == 64
% sols_h = sorted(solve(h, t))
% assert sols_h == [-1, 7]
% # physiquement t >= 0, donc t = 7
% # C5 : h(t) > 3 => -t^2 + 6t + 7 > 3 => -t^2 + 6t + 4 > 0
% # delta = 36 + 16 = 52, racines = 3 +- sqrt(13)
% from sympy import sqrt as Sqrt
% g = -t**2 + 6*t + 4
% delta_g = 36 + 16
% assert delta_g == 52
% sols_g = sorted(solve(g, t))
% t1 = 3 - Sqrt(13)
% t2 = 3 + Sqrt(13)
% assert sols_g[0] == t1
% assert sols_g[1] == t2
% # a < 0, positif entre t1 et t2 (~-0.61 et 6.61), sur [0 ; 7] : ]3-sqrt(13) ; 3+sqrt(13)[
% END-VERIFY
\begin{exercice}{1SPE-SECDEG-EX-030}{2}{20}
Un objet est lancé depuis le sol. Sa hauteur $h$ (en mètres) au bout de $t$ secondes est :
\[h(t) = -t^2 + 6t + 7 \quad (t \geqslant 0).\]

\begin{enumerate}
  \item \textbf{(C3)} Résoudre $h(t) = 0$ dans $\mathbb{R}$. En déduire l'instant auquel l'objet retouche le sol.
  \item Calculer $h(0)$. Interpréter.
  \item \textbf{(C5)} Résoudre l'inéquation $h(t) > 3$, c'est-à-dire déterminer les instants pour lesquels l'objet se trouve à plus de $3$~m du sol.

  \textit{On ramènera à $-t^2 + 6t + 4 > 0$ et on donnera les bornes sous forme exacte.}

  \item \textbf{(C5)} En tenant compte du domaine physique ($0 \leqslant t \leqslant 7$), écrire l'ensemble de solutions.
\end{enumerate}
\end{exercice}
